(1) 对 $n\ne0$，牛顿方程 $m\ddot u_n=K(u_{n+1}+u_{n-1}-2u_n)$；缺陷位置则为 $M\ddot u_0=K(u_1+u_{-1}-2u_0)$。若令 $u_n(t)=U_ne^{-i\omega t}$，得到关于离散振幅 $U_n$ 的二阶差分本征方程。
(2) 在 $n\ge1$ 区域代入 $U_n=A(-1)^n e^{-\kappa n}$，有 $U_{n+1}+U_{n-1}=-2\cosh\kappa\,U_n$。所以 $-m\omega^2U_n=-2K(1+\cosh\kappa)U_n$，即 $\omega^2=(2K/m)(1+\cosh\kappa)>4K/m$，在均匀链带顶之上。
(3) 在缺陷处 $U_0=A,U_{\pm1}=-Ae^{-\kappa}$，因此 $M\omega^2=2K(1+e^{-\kappa})$。结合无缺陷处 $\omega^2=(K/m)(2+e^\kappa+e^{-\kappa})$，约去共同因子后得到 $M/m=2/(1+e^\kappa)$，即 $e^\kappa=2m/M-1$，由于 $M<m$ 所以 $\kappa>0$。
(4) 将 $e^\kappa=(2m-M)/M$ 代回频率，得到 $\omega_{\rm loc}^2=4Km/[M(2m-M)]$。当 $M\to m^-$，$\kappa\to0$、$\omega\to2\sqrt{K/m}$，局域态并入声学带顶；当 $M\to0^+$，频率趋无穷、局域程度增强，展示轻缺陷可形成高频带外振动。
